% Independent proof audit: this fragment has no external document preamble.
\section{An elementary certificate for adaptive four-block contraction}
\label{sec:k4-elementary-certificate}

Let the six edges of the complete graph on four vertices have nonnegative
weights $w_e$. These weights are logarithms of index-domain sizes; opposite
directions have already been aggregated. Suppose a rectangular matrix product
with index logarithms $x,y,z$ costs, up to polynomial factors,
\[
 2^{x+y+z-(3/5)\min\{x,y,z\}}.
\]
For each of the three perfect matchings $M=\{ab,cd\}$, put
$C_M=w_{ac}+w_{ad}+w_{bc}+w_{bd}$. Contract the pair $ab$ and the pair $cd$
independently, then take the inner product of their boundary tables.
The two merge exponents are
\begin{align*}
 C_M+w_{ab}-\tfrac35\min\{w_{ab},w_{ac}+w_{ad},w_{bc}+w_{bd}\},\\
 C_M+w_{cd}-\tfrac35\min\{w_{cd},w_{ac}+w_{bc},w_{ad}+w_{bd}\}.
\end{align*}
Their maximum dominates the final inner-product exponent $C_M$.
Let $F_M(w)$ be this maximum and $F(w)=\min_M F_M(w)$.

\begin{theorem}[Sharp four-block Euclidean certificate]
\label{thm:k4-euclidean}
For every nonnegative $w\in\mathbb R^6$,
\[
 F(w)\le\sqrt{\frac{244}{75}}\,\|w\|_2.
\]
The numerical constant is best possible for these three contraction choices
and this rectangular matrix-product cost model.
\end{theorem}

\begin{proof}
Each $F_M$ is the maximum of six linear forms. For either internal edge
$e\in M$, one form has coefficient vector
\[
 A_e=\mathbf 1_{E\setminus M}+\tfrac25\mathbf 1_e.
\]
The other two forms are obtained by replacing it by
\[
 B_e=A_e+\tfrac35(\mathbf 1_e-\mathbf 1_f-\mathbf 1_g),
\]
where $f,g$ are the two crossing edges incident to one endpoint of $e$.
All coefficient vectors are nonnegative.

Choose one of the six forms from each matching, arbitrarily. Temporarily
replace every selected $B$ form by its corresponding $A$ form. Since every
edge is crossing for exactly two matchings and the three chosen internal
edges are distinct, the sum of the resulting vectors is
\[
 b=2\mathbf 1+\tfrac25 z,
 \qquad z\in\{0,1\}^6,\quad\sum_ez_e=3.
\]
Consequently $\|b\|_2^2=3(12/5)^2+3\cdot2^2=732/25$.

Suppose $k\le3$ selected forms were of type $B$, and let
$D=\sum(\mathbf 1_e-\mathbf 1_f-\mathbf 1_g)$ be the sum of their
correction vectors. We have $\sum_eD_e=-k$. Let $N\ge0$ count, with
multiplicity, occurrences of a negative correction on one of the three
edges selected by $z$. Every positive correction occurs on such an edge,
so $z\mathbin\cdot D=k-N$. It follows that
\[
 b\mathbin\cdot D=-\tfrac85k-\tfrac25N\le-\tfrac85k.
\]
Each individual correction has squared norm three. For corrections from
different matchings, their positive edges are distinct and their negative
two-edge sets intersect in at most one edge. Their inner product is
therefore at most one. Hence
\[
 \|D\|_2^2\le3k+2\binom k2=k(k+2)\le5k.
\]
Restoring the $B$ forms gives a sum vector $b+(3/5)D$, whose norm obeys
\begin{align*}
 \|b+\tfrac35D\|_2^2-\|b\|_2^2
 &=\tfrac65 b\mathbin\cdot D+\tfrac9{25}\|D\|_2^2\\
 &\le-\tfrac{48}{25}k+\tfrac{45}{25}k\le0.
\end{align*}
Thus the mean of any one-form-per-matching selection has squared norm at
most $(732/25)/9=244/75$.

For the given $w$, choose a maximizing form from each matching. The minimum
of the resulting three values is at most their mean. Cauchy--Schwarz and
the preceding uniform estimate prove
$F(w)\le\sqrt{244/75}\,\|w\|_2$.

For sharpness, give the three edges incident to one vertex weight $4/5$
and the other three edges weight $2/3$. Every matching contains one edge
of each weight. Its crossing sum is $44/15$, and its larger merge cost is
$44/15+(2/5)(4/5)=244/75$. The weight vector has squared norm $244/75$,
so equality holds.
\end{proof}

For an application with $r=4b$ original groups divided into four blocks of
size $b$, let $h_{gh}\ge0$ denote directed-pair index logarithms. Each
undirected block-pair weight sums $2b^2$ such logarithms, and the six sets
are disjoint. Therefore
\[
 \|w\|_2^2\le2b^2\sum_{g\ne h}h_{gh}^2.
\]
If the established weighted entropy estimate implies
$\|h\|_2\le(1+o(1))K/r$, the final four-block contraction has coefficient
\[
 \sqrt{\frac{244}{75}}\frac{\sqrt2b}{r}
 =\sqrt{\frac{61}{150}}
\]
on the scale $K$. Formation of the four input block tables, paid recursive
oracle evaluations, storage, packing conditions, and the complete recursion
remain separate obligations. The theorem by itself makes no SAT complexity
claim.

\paragraph{Dependence on the matrix-multiplication exponent.}
The same argument works for every $7/3\le\omega\le3$ when the rectangular
cost is $x+y+z-(3-\omega)\min\{x,y,z\}$. Put $\nu=\omega-2$.
Then $A_e=\mathbf 1_{E\setminus M}+\nu\mathbf 1_e$, and a type-$B$
correction is $(1-\nu)(\mathbf 1_e-\mathbf 1_f-\mathbf 1_g)$.
The base sum is $b=2\mathbf 1+\nu z$, giving
\[
 \|b\|_2^2=3(2+\nu)^2+12=3\omega^2+12,
 \qquad b\mathbin\cdot D\le-(2-\nu)k.
\]
As above $\|D\|_2^2\le5k$, so restoring the corrections changes the
squared norm by at most
\[
 \bigl[-2(1-\nu)(2-\nu)+5(1-\nu)^2\bigr]k
 =(1-\nu)(1-3\nu)k\le0.
\]
Thus the sharp local constant is $\sqrt{(\omega^2+4)/3}$, attained by
three star edges of weight $\omega/3$ and three other edges of weight
$2/3$. The equal-four-cluster entropy coefficient becomes
\[
 \sqrt{\frac{\omega^2+4}{24}}.
\]
This is a parametric contraction theorem; use of a particular $\omega$
still requires an established exact arithmetic implementation of the
corresponding matrix multiplication bound.
